\chapter{Câu hỏi củng cố (Question Sets)}
\label{chap:Question Sets}

\section*{Module 8: The Return and Risk of a Financial Portfolio}
\addcontentsline{toc}{section}{Module 8: The Return and Risk of a Financial Portfolio}

\subsection{Question Set 1}
\label{subsec:app_Question Set 1}

\begin{tcolorbox}[
	colback=gray!5!white,
	colframe=black!15!white,
	boxrule=0.5pt,
	arc=3pt,
	title=\textbf{Question 1},
	coltitle=black,
	colbacktitle=black!5!white,
	before skip=6pt,
	after skip=6pt,
	breakable
]
A portfolio manager creates a portfolio with the following two securities:
\begin{center}
\begin{tabular}{c  c  c}
	\toprule
	\textbf{Security} & \textbf{Weight (\%)} & \textbf{Standard Deviation (\%)} \\
	\midrule
	1 & 30 & 20 \\
	2 & 70 & 12 \\
	\bottomrule
\end{tabular}
\end{center}
If the correlation of returns between the two securities is $0.40$, the expected standard deviation of the portfolio is closest to:
\begin{enumerate}[label=\Alph*.]
	\item 1.5\%.
	\item 8.2\%.
	\item 12.1\%.
\end{enumerate}

\textbf{Solution:}
\textbf{C là phương án đúng.} Áp dụng công thức tính độ lệch chuẩn danh mục hai tài sản:
\[
\sigma_P = \sqrt{w_1^2 \sigma_1^2 + w_2^2 \sigma_2^2 + 2 w_1 w_2 \rho_{1,2} \sigma_1 \sigma_2}
\]
\[
\sigma_P = \sqrt{(0.3^2 \times 0.20^2) + (0.7^2 \times 0.12^2) + 2(0.3)(0.7)(0.40)(0.20)(0.12)}
\]
\[
\sigma_P = \sqrt{0.0036 + 0.007056 + 0.004032} = \sqrt{0.014688} \approx 12.12\%
\]
\textit{Phương án A sai vì đây là phương sai danh mục (1.47\% hoặc 0.0147). Phương án B sai do trừ nhầm phần hiệp phương sai thay vì cộng.}
\end{tcolorbox}

\begin{tcolorbox}[
	colback=gray!5!white,
	colframe=black!15!white,
	boxrule=0.5pt,
	arc=3pt,
	title=\textbf{Question 2},
	coltitle=black,
	colbacktitle=black!5!white,
	before skip=6pt,
	after skip=6pt,
	breakable
]
Using the same weights and standard deviations from Question 1, if the covariance of returns between the two securities is expected to be $-0.0240$ (or $-240\%^2$ in scaled terms), the expected standard deviation of the portfolio is closest to:
\begin{enumerate}[label=\Alph*.]
	\item 0.1\%.
	\item 2.4\%.
	\item 10.2\%.
\end{enumerate}

\textbf{Solution:}
\textbf{B là phương án đúng.} Lưu ý rằng đề bài đã cho trực tiếp hiệp phương sai ($\sigma_{1,2} = -0.0240$) chứ không phải hệ số tương quan.
\[
\sigma_P^2 = w_1^2 \sigma_1^2 + w_2^2 \sigma_2^2 + 2 w_1 w_2 \sigma_{1,2}
\]
\[
\sigma_P^2 = (0.3^2 \times 0.20^2) + (0.7^2 \times 0.12^2) + 2(0.3)(0.7)(-0.0240)
\]
\[
\sigma_P^2 = 0.0036 + 0.007056 - 0.01008 = 0.000576
\]
Độ lệch chuẩn danh mục: $\sigma_P = \sqrt{0.000576} = 2.40\%$.
\textit{Phương án A biểu thị phương sai danh mục (0.0576\% hoặc 0.0006). Phương án C sai do coi nhầm giá trị hiệp phương sai là hệ số tương quan.}
\end{tcolorbox}

\begin{tcolorbox}[
	colback=gray!5!white,
	colframe=black!15!white,
	boxrule=0.5pt,
	arc=3pt,
	title=\textbf{Question 3},
	coltitle=black,
	colbacktitle=black!5!white,
	before skip=6pt,
	after skip=6pt,
	breakable
]
A portfolio of two securities has a return of 15\%. Security 1 has a return of 16\% and Security 2 has a return of 12\%. The proportion invested in Security 1 is:
\begin{enumerate}[label=\Alph*.]
	\item 25\%.
	\item 50\%.
	\item 75\%.
\end{enumerate}

\textbf{Solution:}
\textbf{C là phương án đúng.} Gợi $w_1$ là tỷ trọng đầu tư vào Security 1:
\[
R_P = w_1 R_1 + (1 - w_1) R_2 \implies 15\% = w_1(16\%) + (1 - w_1)(12\%)
\]
\[
15\% = 16\% w_1 + 12\% - 12\% w_1 \implies 3\% = 4\% w_1 \implies w_1 = 0.75 \text{ (hoặc 75\%)}
\]
\end{tcolorbox}

\begin{tcolorbox}[
	colback=gray!5!white,
	colframe=black!15!white,
	boxrule=0.5pt,
	arc=3pt,
	title=\textbf{Question 4},
	coltitle=black,
	colbacktitle=black!5!white,
	before skip=6pt,
	after skip=6pt,
	breakable
]
As the number of assets in an equally weighted portfolio increases, if all assets have the same variance and positive correlation with each other, the contribution of each individual asset's standard deviation to the overall portfolio standard deviation:
\begin{enumerate}[label=\Alph*.]
	\item increases.
	\item decreases.
	\item remains constant.
\end{enumerate}

\textbf{Solution:}
\textbf{B là phương án đúng.} Dựa trên Phương trình~\ref{eq:equal_weighted_variance_m8}, khi số lượng tài sản $N$ tăng lên, cấu phần phương sai riêng lẻ của từng tài sản ($\frac{\sigma^2}{N}$) giảm dần về 0 và trở nên không đáng kể, phản ánh rủi ro phi hệ thống bị đa dạng hóa hoàn toàn. Phương sai danh mục lúc này được quyết định gần như toàn bộ bởi hiệp phương sai trung bình giữa các tài sản.
\end{tcolorbox}

\begin{tcolorbox}[
	colback=gray!5!white,
	colframe=black!15!white,
	boxrule=0.5pt,
	arc=3pt,
	title=\textbf{Question 5},
	coltitle=black,
	colbacktitle=black!5!white,
	before skip=6pt,
	after skip=6pt,
	breakable
]
An analyst projects returns for three assets under three equally likely scenarios:
\begin{center}
\begin{tabular}{c  c  c  c}
	\toprule
	\textbf{Asset} & \textbf{Scenario 1 Return (\%)} & \textbf{Scenario 2 Return (\%)} & \textbf{Scenario 3 Return (\%)} \\
	\midrule
	1 & 12 & 0 & 6 \\
	2 & 12 & 6 & 0 \\
	3 & 0 & 6 & 12 \\
	\bottomrule
\end{tabular}
\end{center}
If the analyst constructs equally weighted two-asset portfolios, which combination has the lowest expected standard deviation?
\begin{enumerate}[label=\Alph*.]
	\item Asset 1 and Asset 2.
	\item Asset 1 and Asset 3.
	\item Asset 2 and Asset 3.
\end{enumerate}

\textbf{Solution:}
\textbf{C là phương án đúng.} Tài sản 2 và Tài sản 3 có tương quan nghịch biến hoàn hảo ($\rho = -1.0$). Một danh mục phân bổ tỷ trọng bằng nhau (50/50) giữa hai tài sản này sẽ mang lại suất sinh lời:
\begin{itemize}
	\item Kịch bản 1: $\frac{12\% + 0\%}{2} = 6\%$
	\item Kịch bản 2: $\frac{6\% + 6\%}{2} = 6\%$
	\item Kịch bản 3: $\frac{0\% + 12\%}{2} = 6\%$
\end{itemize}
Do suất sinh lời của danh mục luôn là một hằng số $6\%$ trong mọi kịch bản, độ lệch chuẩn (đo lường rủi ro biến động) sẽ bằng chính xác 0.
\end{tcolorbox}

\subsection{Question Set 2}
\label{subsec:app_Question Set 2}

\begin{tcolorbox}[
	colback=gray!5!white,
	colframe=black!15!white,
	boxrule=0.5pt,
	arc=3pt,
	title=\textbf{Question 1},
	coltitle=black,
	colbacktitle=black!5!white,
	before skip=6pt,
	after skip=6pt,
	breakable
]
The efficient frontier is the set of asset portfolios that offers the lowest:
\begin{enumerate}[label=\Alph*.]
	\item expected return for a given level of risk.
	\item amount of risk for a given level of return.
	\item average covariance among the assets in the portfolio.
\end{enumerate}

\textbf{Solution:}
\textbf{B là phương án đúng.} Đường biên hiệu quả (efficient frontier) đại diện cho tập hợp các danh mục tối ưu hóa suất sinh lời kỳ vọng ứng với từng mức rủi ro, hoặc tối thiểu hóa rủi ro ứng với từng mức suất sinh lời kỳ vọng.
\end{tcolorbox}

\begin{tcolorbox}[
	colback=gray!5!white,
	colframe=black!15!white,
	boxrule=0.5pt,
	arc=3pt,
	title=\textbf{Question 2},
	coltitle=black,
	colbacktitle=black!5!white,
	before skip=6pt,
	after skip=6pt,
	breakable
]
A portfolio consists of two assets, A and B. The variances and covariance are:
\begin{itemize}
	\item $\sigma_A^2 = 0.0400$
	\item $\sigma_B^2 = 0.0900$
	\item $\sigma_{A,B} = 0.0100$
\end{itemize}
The asset weights that minimize the portfolio variance are closest to:
\begin{enumerate}[label=\Alph*.]
	\item $w_A = 0.273, w_B = 0.727$.
	\item $w_A = 0.500, w_B = 0.500$.
	\item $w_A = 0.727, w_B = 0.273$.
\end{enumerate}

\textbf{Solution:}
\textbf{C là phương án đúng.} Áp dụng công thức tính tỷ trọng danh mục tối thiểu phương sai toàn cầu:
\[
w_A^* = \frac{\sigma_B^2 - \sigma_{A,B}}{\sigma_A^2 + \sigma_B^2 - 2\sigma_{A,B}} = \frac{0.09 - 0.01}{0.04 + 0.09 - 2(0.01)} = \frac{0.08}{0.11} \approx 0.727 \text{ (hoặc 72.7\%)}
\]
Tỷ trọng của tài sản B: $w_B^* = 1 - 0.727 = 0.273$ (hoặc 27.3\%).
\end{tcolorbox}

\begin{tcolorbox}[
	colback=gray!5!white,
	colframe=black!15!white,
	boxrule=0.5pt,
	arc=3pt,
	title=\textbf{Question 3},
	coltitle=black,
	colbacktitle=black!5!white,
	before skip=6pt,
	after skip=6pt,
	breakable
]
The portfolio on the efficient frontier with the lowest standard deviation is best described as the:
\begin{enumerate}[label=\Alph*.]
	\item optimal risky portfolio.
	\item global minimum-variance portfolio.
	\item portfolio with the highest return.
\end{enumerate}

\textbf{Solution:}
\textbf{B là phương án đúng.} Danh mục tối thiểu phương sai toàn cầu (global minimum-variance portfolio) là điểm cực bên trái trên đường biên tối thiểu phương sai, đại diện cho mức rủi ro thấp nhất có thể đạt được từ các tài sản rủi ro.
\end{tcolorbox}

\begin{tcolorbox}[
	colback=gray!5!white,
	colframe=black!15!white,
	boxrule=0.5pt,
	arc=3pt,
	title=\textbf{Question 4},
	coltitle=black,
	colbacktitle=black!5!white,
	before skip=6pt,
	after skip=6pt,
	breakable
]
Which of the following is most appropriate for an investor whose risk tolerance and return preferences are relatively high?
\begin{enumerate}[label=\Alph*.]
	\item The minimum-variance portfolio.
	\item A portfolio near the lower left of the efficient frontier.
	\item A portfolio near the upper right of the efficient frontier.
\end{enumerate}

\textbf{Solution:}
\textbf{C là phương án đúng.} Các nhà đầu tư tìm kiếm suất sinh lời kỳ vọng cao hơn bắt buộc phải chấp nhận mức độ biến động (rủi ro) lớn hơn. Các danh mục này nằm ở góc phía trên bên phải của đường biên hiệu quả.
\end{tcolorbox}

\subsection{Question Set 3}
\label{subsec:app_Question Set 3}

\begin{tcolorbox}[
	colback=gray!5!white,
	colframe=black!15!white,
	boxrule=0.5pt,
	arc=3pt,
	title=\textbf{Question 1},
	coltitle=black,
	colbacktitle=black!5!white,
	before skip=6pt,
	after skip=6pt,
	breakable
]
The line depicting the total risk and expected return of portfolio combinations of a risk-free asset and any risky asset is the:
\begin{enumerate}[label=\Alph*.]
	\item capital market line.
	\item security market line.
	\item capital allocation line.
\end{enumerate}

\textbf{Solution:}
\textbf{C là phương án đúng.} Đường phân bổ vốn (CAL) mô tả các kết hợp giữa tài sản phi rủi ro và một danh mục tài sản rủi ro bất kỳ. CML là một trường hợp đặc biệt của CAL khi danh mục tài sản rủi ro chính là Danh mục thị trường (Market Portfolio). SML mô tả mối quan hệ giữa suất sinh lời kỳ vọng và rủi ro hệ thống (Beta) chứ không phải tổng rủi ro.
\end{tcolorbox}

\begin{tcolorbox}[
	colback=gray!5!white,
	colframe=black!15!white,
	boxrule=0.5pt,
	arc=3pt,
	title=\textbf{Question 2},
	coltitle=black,
	colbacktitle=black!5!white,
	before skip=6pt,
	after skip=6pt,
	breakable
]
The standard deviation of a stock's returns is $30\%$ and the standard deviation of market returns is $25\%$. If the correlation between the stock and market returns is $0.75$, the stock's Beta is closest to:
\begin{enumerate}[label=\Alph*.]
	\item 0.9.
	\item 1.1.
	\item 1.2.
\end{enumerate}

\textbf{Solution:}
\textbf{A là phương án đúng.} Áp dụng công thức tính hệ số Beta của cổ phiếu:
\[
\beta_i = \rho_{i,M} \frac{\sigma_i}{\sigma_M} = 0.75 \left(\frac{30\%}{25\%}\right) = 0.75 \times 1.20 = 0.90
\]
\end{tcolorbox}

\begin{tcolorbox}[
	colback=gray!5!white,
	colframe=black!15!white,
	boxrule=0.5pt,
	arc=3pt,
	title=\textbf{Question 3},
	coltitle=black,
	colbacktitle=black!5!white,
	before skip=6pt,
	after skip=6pt,
	breakable
]
An analyst gathers the following information when the risk-free rate is $3\%$ and the expected market return is $9\%$ (implying a market risk premium of 6\%):
\begin{center}
\begin{tabular}{c  c  c}
	\toprule
	\textbf{Security} & \textbf{Standard Deviation (\%)} & \textbf{Beta} \\
	\midrule
	Security 1 & 25 & 1.50 \\
	Security 2 & 15 & 1.40 \\
	Security 3 & 20 & 1.60 \\
	\bottomrule
\end{tabular}
\end{center}
Based on the Capital Asset Pricing Model, the expected return for Security 1 is:
\begin{enumerate}[label=\Alph*.]
	\item 9.0\%.
	\item 12.0\%.
	\item 13.5\%.
\end{enumerate}

\textbf{Solution:}
\textbf{B là phương án đúng.} Áp dụng mô hình định giá tài sản vốn CAPM:
\[
E[R_1] = R_f + \beta_1 (E[R_M] - R_f) = 3\% + 1.50(9\% - 3\%) = 3\% + 1.50(6\%) = 12\%
\]
\end{tcolorbox}

\begin{tcolorbox}[
	colback=gray!5!white,
	colframe=black!15!white,
	boxrule=0.5pt,
	arc=3pt,
	title=\textbf{Question 4},
	coltitle=black,
	colbacktitle=black!5!white,
	before skip=6pt,
	after skip=6pt,
	breakable
]
Using the table from Question 3, if the risk-free rate is 3\% and the expected market return is 9\%, the market expected return is:
\begin{enumerate}[label=\Alph*.]
	\item 9.0\%.
	\item 12.0\%.
	\item 18.0\%.
\end{enumerate}

\textbf{Solution:}
\textbf{A là phương án đúng.} Suất sinh lời kỳ vọng của thị trường được nêu rõ ràng trong đề bài là $9.0\%$ (Hoặc có thể kiểm chứng lại bằng cách áp dụng CAPM cho danh mục thị trường với $\beta_M = 1.0$: $E[R_M] = 3\% + 1.0(6\%) = 9\%$).
\end{tcolorbox}

\begin{tcolorbox}[
	colback=gray!5!white,
	colframe=black!15!white,
	boxrule=0.5pt,
	arc=3pt,
	title=\textbf{Question 5},
	coltitle=black,
	colbacktitle=black!5!white,
	before skip=6pt,
	after skip=6pt,
	breakable
]
Based on the table in Question 3, a change in the expected market return will have the greatest impact on the expected return of:
\begin{enumerate}[label=\Alph*.]
	\item Security 1.
	\item Security 2.
	\item Security 3.
\end{enumerate}

\textbf{Solution:}
\textbf{C là phương án đúng.} Security 3 có hệ số Beta cao nhất ($\beta = 1.60$), nghĩa là nó nhạy cảm nhất trước các biến động của thị trường chung.
\end{tcolorbox}

\newpage

\section*{Module 9: Simulation of Financial Asset Prices and Returns}
\addcontentsline{toc}{section}{Module 9: Simulation of Financial Asset Prices and Returns}

\subsection{Question Set 1}
\label{subsec:app_mod9_Question_Set_1}

\begin{tcolorbox}[
	colback=gray!5!white,
	colframe=black!15!white,
	boxrule=0.5pt,
	arc=3pt,
	title=\textbf{Question 1},
	coltitle=black,
	colbacktitle=black!5!white,
	before skip=6pt,
	after skip=6pt,
	breakable
]
A limitation of historical simulation is its:
\begin{enumerate}[label=\Alph*.]
	\item inability to analyze highly non-linear exposures.
	\item difficulty in incorporating instruments driven by new factors.
	\item more difficult implementation relative to other simulation approaches.
\end{enumerate}

\textbf{Solution:} \\
\textbf{B là phương án đúng.} Để mô phỏng lịch sử hoạt động hiệu quả cho một công cụ cụ thể, chúng ta cần có đủ dữ liệu lịch sử của các yếu tố chính thúc đẩy biến số đó. Các công cụ mới sẽ không có dữ liệu lịch sử để đối chiếu (ví dụ: cổ phiếu của doanh nghiệp vừa IPO được 2 tuần). \\
Phương án A sai vì mô phỏng lịch sử hoàn toàn có thể áp dụng cho các cấu phần phi tuyến tính. \\
Phương án C sai vì mô phỏng lịch sử thực chất đơn giản và dễ triển khai hơn các phương pháp khác.
\end{tcolorbox}

\begin{tcolorbox}[
	colback=gray!5!white,
	colframe=black!15!white,
	boxrule=0.5pt,
	arc=3pt,
	title=\textbf{Question 2},
	coltitle=black,
	colbacktitle=black!5!white,
	before skip=6pt,
	after skip=6pt,
	breakable
]
One potential criticism of using historical simulation is that it may:
\begin{enumerate}[label=\Alph*.]
	\item be overly sensitive to outliers.
	\item heavily rely on normal distributions to deliver reliable results.
	\item require additional correlation modeling between simulation variables.
\end{enumerate}

\textbf{Solution:} \\
\textbf{A là phương án đúng.} Do phương pháp này dùng trực tiếp dữ liệu quá khứ, kết quả mô phỏng sẽ nhạy cảm và bị ảnh hưởng lớn bởi các điểm dữ liệu dị biệt (outliers) xuất hiện trong chuỗi thời gian quan sát. \\
Phương án B sai vì mô phỏng lịch sử không yêu cầu giả định phân phối chuẩn. \\
Phương án C sai vì tương quan giữa các biến số được bảo toàn tự nhiên trong dữ liệu lịch sử.
\end{tcolorbox}

\begin{tcolorbox}[
	colback=gray!5!white,
	colframe=black!15!white,
	boxrule=0.5pt,
	arc=3pt,
	title=\textbf{Question 3},
	coltitle=black,
	colbacktitle=black!5!white,
	before skip=6pt,
	after skip=6pt,
	breakable
]
Assume a historical sample of asset return data has a positive average return. How would this fact affect the difference between VaR calculated using historical simulation and VaR calculated using the standard parametric approach, conducted according to market best practices?
\begin{enumerate}[label=\Alph*.]
	\item The VaR using historical simulation will be biased to be greater than the VaR using the standard parametric approach.
	\item The VaR using historical simulation will be biased to be less than the VaR using the standard parametric approach.
	\item We cannot infer an answer from the information given.
\end{enumerate}

\textbf{Solution:} \\
\textbf{C là phương án đúng.} Suất sinh lời trung bình dương không cung cấp đủ thông tin để kết luận sự khác biệt giữa hai ước lượng VaR, vì VaR tập trung vào các kịch bản tồi tệ nhất ở phần đuôi trái của phân phối (các khoản lỗ cực đại). Việc xuất hiện một vài kịch bản sinh lời lớn không ảnh hưởng đến giá trị VaR.
\end{tcolorbox}

\subsection{Question Set 2}
\label{subsec:app_mod9_Question_Set_2}

\begin{tcolorbox}[
	colback=gray!5!white,
	colframe=black!15!white,
	boxrule=0.5pt,
	arc=3pt,
	title=\textbf{Question 1},
	coltitle=black,
	colbacktitle=black!5!white,
	before skip=6pt,
	after skip=6pt,
	breakable
]
Bootstrapping scenarios are typically generated by:
\begin{enumerate}[label=\Alph*.]
	\item sampling historical observations with replacement.
	\item drawing from analytical distribution, such as the normal distribution.
	\item assembling a set of scenarios based on actual historical market movements and analyzing in a sequential manner.
\end{enumerate}

\textbf{Solution:} \\
\textbf{A là phương án đúng.} Bootstrapping tạo kịch bản mô phỏng bằng cách lấy mẫu lặp từ tập dữ liệu lịch sử quan sát có thay thế. \\
Phương án B mô tả phương pháp mô phỏng Monte Carlo dựa trên phân phối toán học xác định. \\
Phương án C mô tả phương pháp mô phỏng lịch sử (Historical Simulation) thuần túy.
\end{tcolorbox}

\begin{tcolorbox}[
	colback=gray!5!white,
	colframe=black!15!white,
	boxrule=0.5pt,
	arc=3pt,
	title=\textbf{Question 2},
	coltitle=black,
	colbacktitle=black!5!white,
	before skip=6pt,
	after skip=6pt,
	breakable
]
Will bootstrapping and historical simulation produce similar results when a dataset includes significant outliers?
\begin{enumerate}[label=\Alph*.]
	\item Yes.
	\item No, bootstrapping would be preferable in this case.
	\item No, historical simulation would be preferable in this case.
\end{enumerate}

\textbf{Solution:} \\
\textbf{B là phương án đúng.} Trong trường hợp dữ liệu chứa các outliers dị biệt, bootstrapping có ưu thế hơn nhờ kỹ thuật lấy mẫu lặp lại ngẫu nhiên giúp làm giảm độ nhạy cảm tổng thể của phân phối trước các điểm dị biệt này, so với việc áp dụng chuỗi thời gian cố định trực tiếp như mô phỏng lịch sử.
\end{tcolorbox}

\begin{tcolorbox}[
	colback=gray!5!white,
	colframe=black!15!white,
	boxrule=0.5pt,
	arc=3pt,
	title=\textbf{Question 3},
	coltitle=black,
	colbacktitle=black!5!white,
	before skip=6pt,
	after skip=6pt,
	breakable
]
An implicit assumption in a bootstrapping simulation is that:
\begin{enumerate}[label=\Alph*.]
	\item past data reflect future data.
	\item the sample represents the population.
	\item the distribution characteristics are known.
\end{enumerate}

\textbf{Solution:} \\
\textbf{B là phương án đúng.} Giả định cốt lõi của bootstrapping là mẫu dữ liệu quan sát đại diện tốt cho tổng thể (quần thể), và việc lặp lại mẫu quan sát phản ánh đúng cấu trúc phân phối của tổng thể. \\
Phương án A là giả định của mô phỏng lịch sử. \\
Phương án C là đặc trưng của mô phỏng Monte Carlo dựa trên phân phối đã biết.
\end{tcolorbox}

\begin{tcolorbox}[
	colback=gray!5!white,
	colframe=black!15!white,
	boxrule=0.5pt,
	arc=3pt,
	title=\textbf{Question 4},
	coltitle=black,
	colbacktitle=black!5!white,
	before skip=6pt,
	after skip=6pt,
	breakable
]
An advantage of the bootstrapping approach to simulation is that it:
\begin{enumerate}[label=\Alph*.]
	\item is not overly sensitive to sample quality.
	\item can be used with autocorrelated time-series data.
	\item can generate reliable estimates from limited data.
\end{enumerate}

\textbf{Solution:} \\
\textbf{C là phương án đúng.} Kỹ thuật bootstrapping đặc biệt hữu ích khi xử lý các mẫu dữ liệu nhỏ (dữ liệu hạn chế), giúp tạo ra ước lượng vững chắc về sai số tiêu chuẩn và khoảng tin cậy của tham số. \\
Phương án A sai vì bootstrapping cực kỳ nhạy cảm với chất lượng mẫu. \\
Phương án B sai vì bootstrapping thông thường phá vỡ tính tự tương quan (autocorrelation) của chuỗi thời gian.
\end{tcolorbox}

\subsection{Question Set 3}
\label{subsec:app_mod9_Question_Set_3}

\begin{tcolorbox}[
	colback=gray!5!white,
	colframe=black!15!white,
	boxrule=0.5pt,
	arc=3pt,
	title=\textbf{Question 1},
	coltitle=black,
	colbacktitle=black!5!white,
	before skip=6pt,
	after skip=6pt,
	breakable
]
A strength of Monte Carlo simulation, in comparison to analytical methods for valuing financial assets and liabilities, is the:
\begin{enumerate}[label=\Alph*.]
	\item speed of calculations.
	\item ability to conduct sensitivity analyses.
	\item increased scope of assets and liabilities that can be analyzed.
\end{enumerate}

\textbf{Solution:} \\
\textbf{C là phương án đúng.} Mô phỏng Monte Carlo không yêu cầu công cụ phải có công thức nghiệm đóng (closed-form solution), do đó mở rộng phạm vi định giá cho các tài sản phức tạp, phi tuyến tính như quyền chọn kiểu Mỹ, kiểu Á hay chứng khoán bảo đảm bằng thế chấp (MBS). \\
Phương án A sai vì Monte Carlo rất nặng về tính toán và chậm hơn công thức đóng. \\
Phương án B sai vì cả hai phương pháp đều làm được phân tích độ nhạy.
\end{tcolorbox}

\begin{tcolorbox}[
	colback=gray!5!white,
	colframe=black!15!white,
	boxrule=0.5pt,
	arc=3pt,
	title=\textbf{Question 2},
	coltitle=black,
	colbacktitle=black!5!white,
	before skip=6pt,
	after skip=6pt,
	breakable
]
Relative to analytical methods, a limitation of Monte Carlo simulation is the:
\begin{enumerate}[label=\Alph*.]
	\item inability to perform ``what if'' analysis.
	\item requirement of needing historical data.
	\item lower precision in specifying cause-and-effect relationships.
\end{enumerate}

\textbf{Solution:} \\
\textbf{C là phương án đúng.} Mô phỏng Monte Carlo chỉ cung cấp các ước lượng thống kê xấp xỉ chứ không cho kết quả chính xác tuyệt đối như công thức giải tích. Các công thức giải tích giúp hiểu sâu hơn về mối quan hệ nhân quả giữa các tham số đầu vào và đầu ra.
\end{tcolorbox}

\begin{tcolorbox}[
	colback=gray!5!white,
	colframe=black!15!white,
	boxrule=0.5pt,
	arc=3pt,
	title=\textbf{Question 3},
	coltitle=black,
	colbacktitle=black!5!white,
	before skip=6pt,
	after skip=6pt,
	breakable
]
Which of the following best summarizes the main idea behind using the Cholesky decomposition in a simulation process?
\begin{enumerate}[label=\Alph*.]
	\item Correlated random normal variables can be obtained from simple random normal variables through applying a mathematical technique.
	\item Excess noise from the input data can be eliminated, leading to stronger results.
	\item More efficient simulations can be conducted by being more strategic on how random scenarios are generated.
\end{enumerate}

\textbf{Solution:} \\
\textbf{A là phương án đúng.} Cholesky decomposition là phương pháp toán học giúp biến đổi các biến số ngẫu nhiên độc lập chuẩn hóa thành các biến số ngẫu nhiên chuẩn hóa có tương quan, từ đó mô phỏng các tài sản có mối liên hệ với nhau.
\end{tcolorbox}

\begin{tcolorbox}[
	colback=gray!5!white,
	colframe=black!15!white,
	boxrule=0.5pt,
	arc=3pt,
	title=\textbf{Question 4},
	coltitle=black,
	colbacktitle=black!5!white,
	before skip=6pt,
	after skip=6pt,
	breakable
]
A typical implicit assumption when conducting a Monte Carlo simulation to produce a risk-neutral valuation of a contingent claim is that:
\begin{enumerate}[label=\Alph*.]
	\item the underlying security has zero drift.
	\item each generated path is equally probable.
	\item the calculated value of the contingent claim increases as volatility increases.
\end{enumerate}

\textbf{Solution:} \\
\textbf{B là phương án đúng.} Trong định giá trung lập rủi ro (risk-neutral valuation) bằng Monte Carlo, giả định ngầm định là mỗi đường đi của giá cổ phiếu được mô phỏng có xác suất xảy ra như nhau. Giá trị quyền chọn được tính bằng trung bình cộng các giá trị thanh toán cuối cùng của mọi đường đi rồi chiết khấu về hiện tại.
\end{tcolorbox}
